\begin{center}  \Large{\textbf{ΠΕΡΙΛΗΨΗ}}  \end{center}
\begin{center}  \Large{\textbf{Πρόβλεψη Τιμής και Φορτίου Ηλεκτρικής Ενέργειας\\
στην Ελληνική Αγορά Επόμενης Ημέρας με Μεθόδους Μηχανικής Μάθησης}}  \end{center}
\vspace{3pt}
\vspace{10pt}
\noindent
Η παρούσα διπλωματική εργασία εξετάζει τη βραχυπρόθεσμη πρόβλεψη τιμών της Αγοράς Επόμενης Ημέρας και ηλεκτρικού φορτίου στο ελληνικό σύστημα ηλεκτρικής ενέργειας με χρήση μεθόδων μηχανικής μάθησης. Στόχος είναι η συστηματική αξιολόγηση διαφορετικών στρατηγικών πρόβλεψης για ορίζοντα 24 ωρών και η διερεύνηση της συμπεριφοράς τους σε περιόδους αυξημένης μεταβλητότητας και μεταβολής κατανομής.

Για τον σκοπό αυτό κατασκευάζεται ενιαίο σύνολο δεδομένων από δημόσια διαθέσιμες πηγές, κυρίως την πλατφόρμα ENTSO-E Transparency, εμπλουτισμένο με χρονικές υστερήσεις, ημερολογιακές μεταβλητές και εξωγενείς παράγοντες. Συγκρίνονται τέσσερις στρατηγικές πρόβλεψης, Teacher-Forcing, Recursive, Multi-Input Multi-Output και Direct, σε δύο πλαίσια εκπαίδευσης, στατική εκπαίδευση και μηνιαία walk-forward επανεκπαίδευση. Η αξιολόγηση πραγματοποιείται για πέντε αλγορίθμους, LightGBM, XGBoost, Random Forest, SVR και MLP, με χρήση των δεικτών MAE, RMSE και sMAPE. Επιπλέον, εφαρμόζεται στατιστικός έλεγχος Diebold-Mariano με διόρθωση Harvey-Leybourne-Newbold, ενώ η ανάλυση σημαντικότητας χαρακτηριστικών και ένα διαδραστικό εργαλείο οπτικοποίησης χρησιμοποιούνται για την ερμηνεία και παρουσίαση των αποτελεσμάτων.

Τα αποτελέσματα δείχνουν ότι η μηνιαία walk-forward επανεκπαίδευση αποτελεί κρίσιμο παράγοντα απόδοσης σε εκτεταμένες περιόδους αξιολόγησης. Στο Q1 2026, το καλύτερο retrained LightGBM μοντέλο αναφοράς επιτυγχάνει MAE 10,46 €/MWh στην πρόβλεψη τιμής και 69,4 MW στην πρόβλεψη φορτίου, μειώνοντας το σφάλμα περίπου κατά 28\% και 29\% αντίστοιχα σε σχέση με τη στατική εκπαίδευση. Τα μοντέλα LightGBM και XGBoost εμφανίζουν την πιο σταθερή συνολική συμπεριφορά, ενώ το SVR παρουσιάζει μειωμένη ανθεκτικότητα σε περιόδους μεταβολής κατανομής. Το Teacher-Forcing λειτουργεί ως ιδεατό μοντέλο αναφοράς, χωρίς να είναι άμεσα εφαρμόσιμο στην πράξη, ενώ η Recursive στρατηγική με walk-forward επανεκπαίδευση αναδεικνύεται ως η πιο κατάλληλη επιλογή για επιχειρησιακή χρήση. Η ανάλυση σημαντικότητας χαρακτηριστικών δείχνει ότι η στρατηγική πρόβλεψης επηρεάζει καθοριστικά το είδος της πληροφορίας που αξιοποιείται από τα μοντέλα.

Συνολικά, η εργασία καταλήγει ότι η αξιόπιστη βραχυπρόθεσμη πρόβλεψη στην ελληνική αγορά δεν εξαρτάται αποκλειστικά από την επιλογή αλγορίθμου, αλλά από τον συνδυασμό κατάλληλης στρατηγικής πρόβλεψης, συστηματικής επανεκπαίδευσης και ερμηνεύσιμης αξιολόγησης χωρίς διαρροή πληροφορίας.
\vspace{0.3cm}
\\ \textit{Keywords}: \textbf{electricity price forecasting, load forecasting, Greek Day-Ahead Market, machine learning, forecasting strategies, walk-forward retraining}